\begin{ex}
Tính kết quả của phép tính: $234 + 125$.
\shortans{359}
\loigiai{
Ta có $234 + 125 = 359$.
}
\end{ex}

\begin{ex}
Tính kết quả của phép tính: $342 + 256$.
\shortans{598}
\loigiai{
Ta có $342 + 256 = 598$.
}
\end{ex}

\begin{ex}
Tính kết quả của phép tính: $415 + 372$.
\shortans{787}
\loigiai{
Ta có $415 + 372 = 787$.
}
\end{ex}

\begin{ex}
Tính kết quả của phép tính: $520 + 268$.
\shortans{788}
\loigiai{
Ta có $520 + 268 = 788$.
}
\end{ex}

\begin{ex}
Tính kết quả của phép tính: $631 + 247$.
\shortans{878}
\loigiai{
Ta có $631 + 247 = 878$.
}
\end{ex}

\begin{ex}
Tính kết quả của phép tính: $714 + 185$.
\shortans{899}
\loigiai{
Ta có $714 + 185 = 899$.
}
\end{ex}

\begin{ex}
Tính kết quả của phép tính: $153 + 42$.
\shortans{195}
\loigiai{
Ta có $153 + 42 = 195$.
}
\end{ex}

\begin{ex}
Tính kết quả của phép tính: $324 + 5$.
\shortans{329}
\loigiai{
Ta có $324 + 5 = 329$.
}
\end{ex}

\begin{ex}
Tính kết quả của phép tính: $506 + 72$.
\shortans{578}
\loigiai{
Ta có $506 + 72 = 578$.
}
\end{ex}

\begin{ex}
Tính kết quả của phép tính: $812 + 63$.
\shortans{875}
\loigiai{
Ta có $812 + 63 = 875$.
}
\end{ex}

\begin{ex}
Tính kết quả của phép tính: $400 + 350$.
\shortans{750}
\loigiai{
Ta có $400 + 350 = 750$.
}
\end{ex}

\begin{ex}
Tính kết quả của phép tính: $230 + 145$.
\shortans{375}
\loigiai{
Ta có $230 + 145 = 375$.
}
\end{ex}

\begin{ex}
Tính kết quả của phép tính: $605 + 132$.
\shortans{737}
\loigiai{
Ta có $605 + 132 = 737$.
}
\end{ex}

\begin{ex}
Điền số thích hợp: $241 + \dots = 568$.
\shortans{327}
\loigiai{
Số cần điền: $568 - 241 = 327$.
}
\end{ex}

\begin{ex}
Điền số thích hợp: $352 + \dots = 786$.
\shortans{434}
\loigiai{
Số cần điền: $786 - 352 = 434$.
}
\end{ex}

\begin{ex}
Điền số thích hợp: $\dots + 123 = 457$.
\shortans{334}
\loigiai{
Số cần điền: $457 - 123 = 334$.
}
\end{ex}

\begin{ex}
Điền số thích hợp: $\dots + 240 = 685$.
\shortans{445}
\loigiai{
Số cần điền: $685 - 240 = 445$.
}
\end{ex}

\begin{ex}
Trường Tiểu học Lê Quý Đôn có 324 học sinh nam và 342 học sinh nữ. Trường có tất cả bao nhiêu học sinh?
\shortans{666}
\loigiai{
Có tất cả: $324 + 342 = 666$ học sinh.
}
\end{ex}

\begin{ex}
Một trang trại nuôi 420 con gà và 250 con vịt. Trang trại nuôi tất cả bao nhiêu con gia cầm?
\shortans{670}
\loigiai{
Có tất cả: $420 + 250 = 670$ con.
}
\end{ex}

\begin{ex}
Kho thóc thứ nhất có 315 bao thóc, kho thóc thứ hai có 260 bao thóc. Cả hai kho có bao nhiêu bao thóc?
\shortans{575}
\loigiai{
Cả hai kho có: $315 + 260 = 575$ bao.
}
\end{ex}

\begin{ex}
Bác nông dân thu hoạch được 542 kg bưởi vào đợt một và 325 kg bưởi vào đợt hai. Bác thu hoạch được tất cả bao nhiêu ki-lô-gam bưởi?
\shortans{867}
\loigiai{
Cả hai đợt thu hoạch: $542 + 325 = 867$ kg.
}
\end{ex}

\begin{ex}
Quãng đường từ A đến B dài 235 km, từ B đến C dài 142 km. Quãng đường từ A đến C đi qua B dài bao nhiêu ki-lô-mét?
\shortans{377}
\loigiai{
Dài: $235 + 142 = 377$ km.
}
\end{ex}

\begin{ex}
Tính giá trị biểu thức: $123 + 234 + 341$.
\shortans{698}
\loigiai{
Ta có $123 + 234 = 357$, rồi $357 + 341 = 698$.
}
\end{ex}

\begin{ex}
Tính giá trị biểu thức: $400 + 250 + 130$.
\shortans{780}
\loigiai{
Ta có $400 + 250 = 650$, rồi $650 + 130 = 780$.
}
\end{ex}

\begin{ex}
Tìm tổng của số 432 và số 355.
\shortans{787}
\loigiai{
Tổng: $432 + 355 = 787$.
}
\end{ex}
